\documentclass[11pt,a4paper]{article}
\usepackage[T1]{fontenc}
\usepackage{lmodern,microtype,amsmath,amssymb,amsthm,mathtools}
\usepackage[margin=27mm]{geometry}
\usepackage{booktabs,array,enumitem,xcolor}
\usepackage[colorlinks=true,linkcolor=blue!50!black,citecolor=blue!50!black,urlcolor=blue!50!black]{hyperref}
\hypersetup{pdftitle={Universal optimal tests of unknown unitary dynamics: determinant exponents and tensor stability},pdfauthor={Lluis Eriksson}}
\newtheorem{theorem}{Theorem}[section]
\newtheorem{lemma}[theorem]{Lemma}
\newtheorem{proposition}[theorem]{Proposition}
\newtheorem{corollary}[theorem]{Corollary}
\theoremstyle{remark}\newtheorem{remark}[theorem]{Remark}
\DeclareMathOperator{\Tr}{Tr}
\DeclareMathOperator{\rank}{rank}
\DeclareMathOperator{\Sym}{Sym}
\newcommand{\E}{\mathcal E}
\newcommand{\U}{\mathcal U}
\newcommand{\C}{\mathcal C}
\newcommand{\HH}{\mathcal H}
\newcommand{\EE}{\mathbb E}
\newcommand{\CC}{\mathbb C}
\newcommand{\ket}[1]{\lvert #1\rangle}
\newcommand{\bra}[1]{\langle #1\rvert}
\newcommand{\vket}[1]{\lvert #1\rangle\!\rangle}
\newcommand{\vbra}[1]{\langle\!\langle #1\rvert}
\newcommand{\HS}{\mathrm{HS}}
\newcommand{\bal}{\lambda_*}
\setlength{\emergencystretch}{2em}
\setlist{nosep,leftmargin=*}
\title{Universal optimal tests of unknown unitary dynamics:\\determinant exponents and tensor stability}
\author{Lluis Eriksson\\\small Independent researcher, Sweden}
\date{16 October 2026\quad---\quad Research preprint}
\begin{document}
\maketitle
\begin{abstract}
We study forward-query tests of an arbitrary fixed unitary channel against memoryless noise with unknown fixed input and output rotations. A Schur--Weyl reduction gives a family of optimal candidate tests, but leaves a channel-dependent minimization over partitions. We remove that minimization: the balanced partition is optimal simultaneously for every channel, every query number and every compact alternative family, even against adaptive and more general positive physical testers. The key identity expresses each block acceptance probability as a normalized Schur polynomial averaged over a Gaussian Kraus matrix. Classical majorization then identifies a universal optimizer. The same identity determines the exact asymptotic error exponent through the maximum determinant of a Kraus combination, with explicit polynomial finite-query bounds. Perfect detection is possible precisely when the Kraus span contains no invertible matrix; its first possible query number is one more than the maximum matrix rank in that span. For qubits, all even-query errors are determined by the spin-flip spectrum of the normalized Choi state. We characterize stability of perfect detection under every tensor partner by a shrunk-subspace condition, using classical noncommutative-rank theory. A completely depolarizing partner of dimension $d-1$ destroys perfect detection whenever this condition fails. For the qutrit Werner--Holevo channel we compute the exact exponent after adjoining the smallest possible partner, a depolarizing qubit, and also for its tensor square. The hypotheses, particularly the independent unknown rotations and the global-unitary null in the tensor example, are essential.
\end{abstract}
\noindent\textbf{Keywords:} quantum hypothesis testing; Schur polynomials; Gaussian moments; Kraus matrix spaces; error exponents; noncommutative rank.

\section{Introduction and contribution}
Can one choose a single experiment that optimally detects any departure from an unknown unitary dynamics, without knowing the orientation or the type of noise? In general channel discrimination this is too much to expect. In the rotation-invariant statistical model considered here, however, the answer is affirmative at every finite query number.

The unitary null is the same unknown channel at every call. The alternative is the same memoryless CPTP map at every call, with independent unknown input and output rotations that are fixed throughout the experiment. Performance is measured in the worst case over these rotations. The experiment we identify depends only on the dimension, the query number and the permitted false-positive probability. It achieves the optimal failure probability for each alternative separately, and therefore also for arbitrary compact families of alternatives.

The main mechanism connects quantum testing to classical symmetric-function inequalities. A Gaussian linear combination of Kraus matrices converts the acceptance probability in a Schur--Weyl block into an expectation of a normalized Schur polynomial. The balanced partition minimizes these polynomials pointwise. In that block the Schur polynomial reduces to a determinant power times an elementary symmetric polynomial. This gives an exact error exponent, a finite-query perfect-detection criterion, and explicit qubit formulas.

\paragraph{What builds on earlier work.}
Quantum testers and symmetry-based minimax reductions are established tools~\cite{nakahira}. Multiplicity-entangled probes, balanced dimension/multiplicity ratios, and finite Schmidt-rank truncation appear in Chiribella and Ebler~\cite{chiribella}, notably their Appendix~B. Antisymmetric probes have also been used for unambiguous unitary comparison~\cite{sedlak}. Parallel optimality for compact-group estimation and subgroup testing has substantial precedents~\cite{hayashi,hayashi_symmetry}. These constructions and principles are not claimed as new.

The author's preceding manuscript~\cite{complete} establishes a complete class for the present two-sided rotation model, reducing its optimum to a minimum over partitions. It determines the optimizer for two queries and for complete depolarization, but leaves a general channel-dependent optimization. We reproduce the needed reduction below and remove that optimization for all CPTP maps. The resulting determinant exponent and exact rank threshold are additional conclusions, not consequences of a numerical search.

The majorization inequality itself is classical: Sra~\cite{sra} proves monotonicity of normalized Schur functions under dominance. For the particular balanced partition needed here, a short Muirhead argument suffices. The spin-flip roots used for qubits are the roots in Wootters' concurrence formula~\cite{wootters}; determinant-based entanglement quantities also occur in Gour's work~\cite{gour}. Our claim is their operational role in this specific minimax testing problem, not a new symmetric-function inequality or a new concurrence construction.

Werner--Holevo channels are standard examples~\cite{werner}. Ghosal et al.~\cite{ghosal} recently study testing non-mixed-unitarity, including odd-dimensional Werner--Holevo examples and infinite Stein exponents. Their null includes mixed-unitary channels, their orientation promise differs, and their asymptotic error criterion is not the finite perfect-detection problem here. We make no claim to have first observed the relevance of their parity. Singular matrix spaces and their behavior under tensor constructions are likewise classical algebraic objects~\cite{ggow,derksen}. The explicit testing consequences below must be read with these distinctions.


An additional question is whether a finite perfect test survives enlarging the system and its unknown-unitary null. We give a complete answer through shrunk subspaces of the Kraus span. This is an operational application of classical noncommutative-rank and blow-up theorems, not a new algebraic characterization. It distinguishes channels whose perfect detection survives every tensor partner from those for which even a completely depolarizing partner destroys it. The exact six-dimensional example below quantifies that change and attains the smallest possible partner dimension.

\section{The testing problem and its universal answer}
Fix $d\ge2$ and $n\ge1$. Let $\E$ be a CPTP map on $\CC^d$. Write $\U_W(X)=WXW^\dagger$. The hypotheses are
\begin{align}
H_0 &: \U_W,\qquad W\in U(d),\\
H_1 &: \U_L\circ\E\circ\U_R,\qquad L,R\in U(d).
\end{align}
The same $W$, or the same triple $(\E,L,R)$, is used at all calls. Under $H_1$ the calls are memoryless. There is no access to the environment, inverses, or a controlled implementation of an unspecified dilation. Ancillas, arbitrary causal control and a final binary measurement are allowed.

We use unnormalized Choi matrices,
\[
J_\E=\sum_{i,j}\E(\ket i\bra j)\otimes\ket i\bra j,
\qquad \Tr_OJ_\E=I_I.
\]
All output legs precede all input legs; a tensor power of a Choi matrix includes this fixed regrouping. A physical accepting tester effect $T\succeq0$ has acceptance probability $a_T(\E)=\Tr(TJ_\E^{\otimes n})$. The accepting outcome accepts $H_0$. Define
\[
\alpha(T)=1-\inf_W a_T(\U_W),\qquad
\beta_\E(T)=\sup_{L,R}a_T(\U_L\circ\E\circ\U_R).
\]
Let $\beta_{n,\E}(\alpha)$ be the optimum at $\alpha(T)\le\alpha$. For a family $\C$, also take the supremum over $\E\in\C$. The converse applies to any physically normalized positive-tester class containing parallel protocols and having the displayed Choi representation. It does not identify arbitrary positive matrices with physical experiments.

For partitions $\lambda\vdash n$ with at most $d$ rows, Schur--Weyl duality gives
\[
\HH_n=(\CC^d)^{\otimes n}
=\bigoplus_\lambda V_\lambda\otimes M_\lambda,
\qquad q_\lambda=\dim V_\lambda=s_\lambda(1^d),\quad
m_\lambda=\dim M_\lambda.
\]
Here $s_\lambda$ is the Schur polynomial. Write $n=ad+r$, $0\le r<d$, and
\[
\bal=((a+1)^r,a^{d-r}),\qquad q_*=\binom dr.
\]
Choose any Kraus representation $\E(X)=\sum_{j=1}^kK_jXK_j^\dagger$. Let $g_j$ be independent standard circular complex Gaussian variables, with $\EE|g_j|^2=1$, and put $G=\sum_jg_jK_j$. For a positive matrix, $e_r$ denotes the elementary symmetric polynomial of its eigenvalues; $e_0=1$.

\begin{theorem}[One optimal test for every noise channel]\label{thm:universal}
For every $0\le\alpha\le1$,
\begin{equation}\boxed{\displaystyle
\beta_{n,\E}(\alpha)=(1-\alpha)b_n(\E),\qquad
b_n(\E)=\frac{\EE\bigl[|\det G|^{2a}e_r(G^\dagger G)\bigr]}
{n!\binom dr}.}\label{eq:main}
\end{equation}
At $a=0$ the determinant factor is $1$, including for singular $G$. The same parallel test attains this optimum for all $\E$: prepare
\[
\frac{I_{V_*}}{q_*}\otimes\ket{\phi_{m_*}}\bra{\phi_{m_*}},
\qquad \ket{\phi_m}=m^{-1/2}\sum_{j=1}^m\ket{j,j},
\]
where one multiplicity register is a bypass reference. After the $n$ parallel calls, accept upon return to the same isotypic block and the same maximally entangled multiplicity state, then retain acceptance with probability $1-\alpha$.

For any nonempty compact family $\C$ of CPTP maps,
\begin{equation}
\beta_{n,\C}(\alpha)=(1-\alpha)\max_{\E\in\C}b_n(\E).
\label{eq:robust}
\end{equation}
Adaptation and any larger physical positive-tester class as specified above give the same optimum.
\end{theorem}

The probe and measurement do not depend on $\E$ or on $\C$. Their reference dimension $m_*$ can grow exponentially with $n$; the theorem does not claim optimality with a smaller reference. At $n\le d$, $\bal=(1^n)$ and $m_*=1$: the optimal input is the normalized antisymmetric projector and the output test is that projector, without a bypass reference. At $n=1$, $b_1=1$ for every channel.

\section{Proof: positive testers, Gaussian moments and majorization}
\subsection{The complete-class reduction}
Let $P_\pi$ permute the $n$ physical factors and let $\Pi_\lambda$ project onto an isotypic block. Define
\[
S_n=\frac1{n!}\sum_{\pi\in\mathfrak S_n}P_\pi^O\otimes P_\pi^I,
\qquad D_\lambda=(\Pi_\lambda^O\otimes\Pi_\lambda^I)S_n.
\]
Permutation matrices are real in the computational basis; the input Schur spaces are the conjugate spaces appropriate to vectorization. In compatible Schur bases,
\[
D_\lambda=I_{V_\lambda^O}\otimes I_{\overline V_\lambda^I}
\otimes\ket{\phi_{m_\lambda}}\bra{\phi_{m_\lambda}}.
\]
These are orthogonal projectors of rank $q_\lambda^2$ and sum to $S_n$. Define
\[
b_\lambda(\E)=\frac{\Tr(D_\lambda J_\E^{\otimes n})}{q_\lambda m_\lambda}.
\]

\begin{lemma}[Block reduction, following \cite{complete}]\label{lem:block}
For each physical $T\succeq0$, there are weights $w_\lambda\ge0$ such that
\begin{align}
\sum_\lambda w_\lambda&=c:=\int a_T(\U_W)\,dW\le1,\\
\int a_T(\U_L\circ\E\circ\U_R)\,dL\,dR
&\ge\sum_\lambda w_\lambda b_\lambda(\E).
\end{align}
Each $b_\lambda$ is the acceptance probability of a physical parallel test with null acceptance one. The same weights work for every CPTP map.
\end{lemma}
\begin{proof}
The Haar null moment is
\[
M_0=\int J_{\U_W}^{\otimes n}\,dW
=\sum_\lambda\frac{m_\lambda}{q_\lambda}D_\lambda.
\]
Twirl $J_\E^{\otimes n}$ independently by collective output and input unitaries, obtaining $B_\E$. Schur's lemma makes $B_\E$ the identity on each pair of representation factors, tensored with an operator on their multiplicity factors. The original tensor power, and hence the twirl, commutes with simultaneous permutations. In $M_\lambda\otimes\overline M_\mu$, the trivial permutation representation occurs only when $\lambda=\mu$, with multiplicity one and vector $\phi_{m_\lambda}$. Thus every $D_\lambda$ reduces $B_\E$, and
\[
B_\E=\sum_\lambda\frac{m_\lambda}{q_\lambda}b_\lambda(\E)D_\lambda+B_\E^\perp,
\qquad B_\E^\perp\succeq0.
\]
Set $w_\lambda=(m_\lambda/q_\lambda)\Tr(TD_\lambda)$. Positivity gives the inequality and $\sum w_\lambda=\Tr(TM_0)=c$.

To realize a block, prepare $I_{V_\lambda}/q_\lambda$ and a maximally entangled pair between $M_\lambda$ and a reference, and project the output multiplicity and reference back onto that pair, accepting only in block $\lambda$. Its accepting tester is $D_\lambda/(q_\lambda m_\lambda)$. Collective unitaries act trivially on multiplicity, giving null acceptance one. The accepting tester is invariant under both collective rotations. This verifies its probability and its physical normalization.
\end{proof}

The preceding lemma yields $(1-\alpha)\min_\lambda b_\lambda$ for a single channel, since $c\ge1-\alpha$. The next identity identifies its minimum simultaneously for all channels.

\subsection{A Gaussian identity}
\begin{lemma}[Gaussian--Schur representation]\label{lem:gauss}
For every partition in the decomposition,
\begin{equation}
b_\lambda(\E)=\frac{\EE s_\lambda(G^\dagger G)}{n!s_\lambda(1^d)}.
\label{eq:gauss}
\end{equation}
\end{lemma}
\begin{proof}
Write $v=\vket G=\sum_jg_j\vket{K_j}$, so $\EE vv^\dagger=J_\E$. The complex Gaussian moment formula, obtained by pairing each $g$ with a conjugate $g$, gives
\[
\EE\bigl[(vv^\dagger)^{\otimes n}\bigr]
=n!S_nJ_\E^{\otimes n}S_n.
\]
The permutation symmetry makes the two expressions with one or two $S_n$ equal. Since $D_\lambda\le S_n$,
\[
\Tr(D_\lambda J_\E^{\otimes n})
=\frac1{n!}\EE\|D_\lambda\vket{G^{\otimes n}}\|^2.
\]
For invertible $G$, its Schur decomposition is $G^{\otimes n}=\bigoplus_\lambda R_\lambda(G)\otimes I_{M_\lambda}$. Consequently the squared norm on the right is
\[
m_\lambda\Tr\bigl(R_\lambda(G)^\dagger R_\lambda(G)\bigr)
=m_\lambda s_\lambda(G^\dagger G).
\]
The same formula holds for singular $G$ by polynomial continuity. Division by $q_\lambda m_\lambda$ proves the identity. It also shows that the Gaussian law, and hence the expression, is independent of the chosen Kraus representation.
\end{proof}

\subsection{The balanced partition is always minimal}
\begin{lemma}[Balanced normalized Schur minimum]\label{lem:balanced}
For $x\in[0,\infty)^d$ and every $\lambda\vdash n$ with at most $d$ rows,
\begin{equation}
\frac{s_\lambda(x)}{s_\lambda(1^d)}
\ge\frac{s_{\bal}(x)}{s_{\bal}(1^d)}
=\frac{(x_1\cdots x_d)^a e_r(x)}{\binom dr}.
\end{equation}
\end{lemma}
\begin{proof}
This follows from Sra's dominance theorem~\cite{sra}, since every partition of $n$ into at most $d$ parts dominates $\bal$. Here is a direct proof of the needed special case. Expand $s_\lambda=\sum_\mu K_{\lambda\mu}m_\mu$ in monomial symmetric polynomials. The Kostka coefficients are nonnegative integers. After normalizing at $1^d$, the result is a convex combination of $m_\mu(x)/m_\mu(1^d)$. Every exponent partition $\mu$ has total $n$ and dominates $\bal$. Muirhead's inequality bounds each normalized monomial mean below by that of $\bal$. Finally $s_{\bal}=m_{\bal}=(x_1\cdots x_d)^ae_r(x)$, since the representation is a determinant power times an exterior power. The assertion follows, including zero coordinates by continuity.
\end{proof}

Taking expectations in Lemma~\ref{lem:balanced} and using Lemma~\ref{lem:gauss} gives $b_\lambda(\E)\ge b_{\bal}(\E)$ for every channel. Lemma~\ref{lem:block} now implies
\[
\beta_\E(T)\ge\sum_\lambda w_\lambda b_\lambda(\E)
\ge c\,b_{\bal}(\E)\ge(1-\alpha)b_{\bal}(\E).
\]
The balanced parallel test attains equality. The same inequalities hold for each $\E$ before taking a supremum, proving both parts of Theorem~\ref{thm:universal}. Notice that domination compares with the original tester's orbit average, not its error at every specified orientation.

\section{Determinant geometry gives the exact error exponent}
Let $k=\rank J_\E$ and take a minimal Kraus representation. For $z\in\CC^k$, put $K(z)=\sum_jz_jK_j$, and define
\begin{equation}
\delta(\E)=\max_{\|z\|_2=1}|\det K(z)|^{2/d}.
\label{eq:delta}
\end{equation}
This number does not depend on the minimal Kraus basis. Equivalently,
\begin{equation}
\delta(\E)=\max\{\,|\det A|^{2/d}:\vket A\vbra A\preceq J_\E\,\}.
\label{eq:intrinsic}
\end{equation}
Indeed, for the matrix $V$ whose columns are $\vket{K_j}$, $xx^\dagger\preceq VV^\dagger$ holds exactly when $x=Vz$ for some $\|z\|\le1$. Homogeneity permits replacing the ball by its sphere. This also allows redundant Kraus representations with zero coefficients.

\begin{theorem}[Finite bounds and exact exponent]\label{thm:exponent}
If $\delta(\E)>0$, then, for $n=ad+r$,
\begin{equation}\boxed{\displaystyle
\delta(\E)^n\le b_n(\E)
\le\binom{n+k-1}{k-1}\delta(\E)^{n-r}.}\label{eq:bounds}
\end{equation}
In particular, at multiples of $d$ the two bounds differ only by the displayed polynomial. For fixed $d$ and $0\le\alpha<1$,
\begin{equation}
\lim_{n\to\infty}-\frac1n\log\beta_{n,\E}(\alpha)
=-\log\delta(\E).
\label{eq:exponent}
\end{equation}
Moreover $0\le\delta\le1$, and $\delta=1$ exactly for unitary channels. If $\delta=0$, $b_n=0$ for every $n\ge d$.
\end{theorem}

\begin{proof}
First $K(z)$ is a contraction for $\|z\|=1$: Cauchy--Schwarz gives
\[
\|K(z)v\|^2\le\sum_j\|K_jv\|^2=\|v\|^2.
\]
Let $\sigma$ be normalized surface measure on the complex unit sphere and $N_{n,k}=\binom{n+k-1}{k-1}$. The radial decomposition $g=Rz$ of a complex Gaussian gives $\EE R^{2n}=n!N_{n,k}$. Formula~\eqref{eq:main} becomes
\begin{equation}
b_n=\frac{N_{n,k}}{\binom dr}\int
|\det K(z)|^{2a}e_r(K(z)^\dagger K(z))\,d\sigma(z).
\label{eq:sphere}
\end{equation}
The contraction property bounds $e_r/\binom dr\le1$ and gives the upper bound in~\eqref{eq:bounds}.

For the lower bound, consider the vector-valued homogeneous polynomial of degree $n$
\[
F(z)=\det(K(z))^a\,\bigwedge\nolimits^r K(z).
\]
For $r=0$, the exterior power is the scalar $1$. Its squared Hilbert--Schmidt norm is the integrand of~\eqref{eq:sphere}. For every degree-$n$ homogeneous polynomial with values in a finite-dimensional Hilbert space,
\begin{equation}
\|F(z)\|^2\le N_{n,k}\int\|F(u)\|^2\,d\sigma(u),\qquad\|z\|=1.
\label{eq:kernel}
\end{equation}
To see this directly, use the orthogonal monomials $z^\nu$, $|\nu|=n$, whose sphere norms squared are $\nu!/(n!N_{n,k})$. Cauchy--Schwarz and the multinomial theorem give the evaluation norm $N_{n,k}$. Applying the scalar inequality componentwise proves the vector case.

Choose a maximizing $z$ in~\eqref{eq:delta}. Maclaurin's inequality gives $e_r(K^\dagger K)/\binom dr\ge|\det K|^{2r/d}$. Combining this with~\eqref{eq:kernel} proves $b_n\ge\delta^n$. Taking logarithms proves~\eqref{eq:exponent}.

If $\delta=1$, the maximizing contraction $A$ has every singular value one, so is unitary. Equation~\eqref{eq:intrinsic} implies $J_\E-\vket A\vbra A\succeq0$, with zero partial trace; hence it is zero and $\E=\U_A$. The converse is immediate. Finally, $\delta=0$ means the determinant polynomial vanishes identically, so~\eqref{eq:main} is zero whenever $a\ge1$.
\end{proof}

\begin{corollary}[Robust exponent]\label{cor:robust-exp}
Fix $0\le\alpha<1$. For a nonempty compact family $\C$, let $\delta_\C=\max_{\E\in\C}\delta(\E)$. If $\delta_\C>0$, the worst-case exponent is $-\log\delta_\C$, with the uniform bounds~\eqref{eq:bounds} using $k=d^2$ and $\delta=\delta_\C$. If $\delta_\C=0$, the universal test is perfect for every member at $n=d$. In particular, every compact family containing no unitary channel has a strictly positive exponent, or finite perfect detection.
\end{corollary}
\begin{proof}
Use a $d^2$-column square root of the Choi matrix to parameterize~\eqref{eq:intrinsic}. Continuity of the positive square root and maximization over a fixed compact ball show that $\delta$ is continuous. Take suprema in~\eqref{eq:bounds}, use~\eqref{eq:robust}, and apply Theorem~\ref{thm:exponent}.
\end{proof}

The determinant maximum can be a difficult polynomial optimization. The theorem gives an exact invariant and bounds; it does not provide an efficient general algorithm. Nor is $\delta$ the maximum weight of a unitary channel below $\E$ in the CP order. For nontrivial qubit amplitude damping that weight is zero, while $\delta$ is positive, as Section~\ref{sec:examples} shows.

\section{An exact criterion and threshold for perfect detection}
Let
\[
\rho(\E)=\max\{\rank A:A\in\operatorname{span}_{\CC}(K_1,\ldots,K_k)\}.
\]
This is ordinary matrix-space rank over commuting coefficients. It is independent of Kraus representation and is not noncommutative rank.

\begin{theorem}[First perfect query number]\label{thm:perfect}
For $0\le\alpha<1$, a finite-query test with $\beta_{n,\E}(\alpha)=0$ exists if and only if $\rho(\E)<d$. In that case its smallest query number, for all the strategy classes in Theorem~\ref{thm:universal}, is
\begin{equation}
n_{\mathrm{perfect}}=\rho(\E)+1.
\end{equation}
The threshold is attained with an antisymmetric input and output projector, without a bypass reference.
\end{theorem}
\begin{proof}
For $n<d$, formula~\eqref{eq:main} reduces to $\EE e_n(G^\dagger G)/(n!\binom dn)$. Its integrand is a sum of squared $n\times n$ minors. The expectation is zero exactly when all these polynomials vanish identically, equivalently when every Kraus combination has rank below $n$. A nonzero polynomial has positive squared Gaussian integral. At $n=d$ the same argument uses the determinant. If an invertible combination exists, $\delta>0$ and Theorem~\ref{thm:exponent} rules out zero at every $n$. If none exists, zero is attained at $\rho+1\le d$, while all smaller $n$ have strictly positive error. The balanced block at that threshold is antisymmetric with multiplicity one.
\end{proof}

Thus a question about all adaptive discrimination experiments becomes a question about the largest nonvanishing minor of a matrix pencil. This is a structural equivalence, not a claim of deterministic polynomial-time solution of general symbolic determinant identity testing.


\section{Which perfect tests survive every tensor partner?}\label{sec:stability}
Let $S=\operatorname{span}\{K_j\}\subseteq M_d(\CC)$ be the Kraus matrix space. For $U\subseteq\CC^d$, write $S(U)=\operatorname{span}\{Au:A\in S,u\in U\}$, and define
\[
c(S)=\max_{U\subseteq\CC^d}\bigl(\dim U-\dim S(U)\bigr).
\]
The zero subspace gives $c(S)\ge0$. A subspace with positive deficiency is called a \emph{shrunk subspace}. For a square channel $\mathcal F$ on $\CC^m$, testing $\E\otimes\mathcal F$ means the original model in dimension $dm$: the null contains every fixed $U(dm)$ unitary, and the alternative has independent unknown global input and output rotations. This is not the task obtained by adding a freely bypassing reference or retaining a product-unitary null.

\begin{theorem}[Tensor stability via a classical matrix-space criterion]\label{thm:stable}
Fix $0\le\alpha<1$. The following are equivalent:
\begin{enumerate}
\item For every $m\ge1$ and every CPTP map $\mathcal F$ on $\CC^m$, the channel $\E\otimes\mathcal F$ admits a finite perfect test in the global model.
\item $S$ has a shrunk subspace, or equivalently $c(S)>0$.
\item The noncommutative rank of $S$ is less than $d$.
\end{enumerate}
If $c=c(S)>0$, then for every such partner
\begin{equation}
\rho(S\otimes S_{\mathcal F})\le m(d-c),\qquad
n_{\mathrm{perfect}}(\E\otimes\mathcal F)\le m(d-c)+1.
\label{eq:stable-bound}
\end{equation}
If $c=0$, the completely depolarizing partner $\mathcal D_{d-1}$ already gives
\begin{equation}
\delta(\E\otimes\mathcal D_{d-1})>0,
\label{eq:unstable}
\end{equation}
so no finite perfect test exists for that partner. Here $\mathcal D_1$ is the one-dimensional identity channel.
\end{theorem}
\begin{proof}
If $U$ has deficiency $c$, every matrix in $S\otimes S_{\mathcal F}$ maps $U\otimes\CC^m$ into $S(U)\otimes\CC^m$. Its kernel therefore has dimension at least $mc$. This proves the rank bound, and Theorem~\ref{thm:perfect} gives the query bound. The resulting antisymmetric test requires no bypass reference since its query number is at most $dm$.

For completeness, the classical noncommutative-rank characterization says
\[
\operatorname{ncrk}(S)=d-c(S),
\]
and $\operatorname{ncrk}(S)=d$ if and only if a tensor blow-up $S\otimes M_m(\CC)$ contains an invertible matrix. These are established results of matrix-space theory; see the attributions and equivalent formulations in Ivanyos--Qiao--Subrahmanyam~\cite{iqs} and Li et al.~\cite{li}, Section II.2 and Theorem 9. Derksen--Makam~\cite{degree}, Theorem 1.8, further shows that when the rank is full an invertible element exists in every blow-up size $m\ge d-1$. Since the Kraus span of $\mathcal D_m$ is exactly $M_m(\CC)$, taking $m=d-1$ gives an invertible Kraus combination for $\E\otimes\mathcal D_{d-1}$. Theorems~\ref{thm:exponent} and~\ref{thm:perfect} prove~\eqref{eq:unstable} and rule out condition 1. The same classical identity proves equivalence with condition 3.
\end{proof}

\begin{corollary}[Unital perfect detection is fragile]
If $\E$ is unital, then $c(S)=0$. In particular, whenever a unital channel admits a finite perfect test, some depolarizing partner of dimension at most $d-1$ removes that possibility in the global model.
\end{corollary}
\begin{proof}
For the orthogonal projection $P_U$, positivity and unitality give $0\preceq\E(P_U)\preceq I$. Its support is $S(U)$ and trace preservation gives $\Tr\E(P_U)=\dim U$. Hence $\dim U\le\dim S(U)$ for all $U$.
\end{proof}

The distinction has an algorithmic consequence. When the Kraus span is supplied by a rational matrix basis, deterministic polynomial-time noncommutative-rank algorithms~\cite{ggow,iqs} decide tensor stability. This invokes those algorithms rather than solving commutative symbolic determinant identity testing. It also does not supply an efficient optimizer for $\delta$ or an efficient physical implementation of a large-query test.

The stable class is nonempty: for a pure replacer on $\CC^d$, $S(\CC^d)$ is one-dimensional, so $c=d-1$. Equation~\eqref{eq:stable-bound} yields $n_{\mathrm{perfect}}\le m+1$ for every $m$-dimensional partner. For the partner $\mathcal D_m$, the matrix space consists of all maps into a fixed $m$-dimensional output space and has maximum rank $m$, so this bound is attained.

\section{Qubit channels: the spin-flip spectrum determines every even error}
Let $d=2$ and choose a minimal Kraus representation, so $k\le4$. Define the complex symmetric matrix $C$ by
\[
\det K(z)=z^TCz,\quad C_{jj}=\det K_j,\quad
C_{ij}=\tfrac12\bigl(\det(K_i+K_j)-\det K_i-\det K_j\bigr)\quad(i\ne j).
\]
Let $\sigma_1\ge\cdots\ge\sigma_k\ge0$ be its Takagi singular values, padded with zeros to four values when necessary.

\begin{theorem}[Exact qubit formulas]\label{thm:qubit}
For $a\ge1$,
\begin{align}
\delta(\E)&=\sigma_1,\\
b_{2a}(\E)&=\frac{(a!)^2}{(2a)!}[t^a]
\prod_{j=1}^k(1-4\sigma_j^2t)^{-1/2}.
\label{eq:qubit}
\end{align}
If $\omega=J_\E/2$ and $\widetilde\omega=(Y\otimes Y)\overline\omega(Y\otimes Y)$, where $Y$ is Pauli $Y$, then these $\sigma_j$ are the eigenvalues of
\[
\sqrt{\sqrt\omega\,\widetilde\omega\,\sqrt\omega}.
\]
They are the spin-flip roots entering Wootters' concurrence formula, although here the exponent depends on their largest value, not on the concurrence combination.
\end{theorem}
\begin{proof}
Takagi factorization and unitary invariance of the complex Gaussian reduce $\det G$ to $\sum_j\sigma_jg_j^2$. Its maximum absolute value on the unit sphere is $\sigma_1$. Expanding its $a$th power and pairing the independent phases gives
\[
\EE\left|\sum_j\sigma_jg_j^2\right|^{2a}
=(a!)^2\sum_{a_1+\cdots+a_k=a}\prod_j
\frac{(2a_j)!}{(a_j!)^2}\sigma_j^{2a_j}.
\]
The generating series of the central binomial coefficients is $(1-4t)^{-1/2}$, proving~\eqref{eq:qubit}.

For the spectral identification, set $v_j=\vket{K_j}/\sqrt2$, so $\omega=\sum_jv_jv_j^\dagger$. The symmetric preconcurrence matrix $\tau_{ij}=v_i^\dagger(Y\otimes Y)\overline v_j$ equals $-\overline C_{ij}$ in this vectorization convention, by the polarized determinant identity. Its singular values are therefore those of $C$. The standard ensemble identity in Wootters' proof~\cite{wootters}, or the equality of nonzero spectra of rectangular matrix products, identifies these with the displayed spin-flip roots.
\end{proof}

For example, $b_2=\sum_j\sigma_j^2$, while
\[
b_4=\sum_j\sigma_j^4+\frac23\sum_{i<j}\sigma_i^2\sigma_j^2.
\]
These formulas require no minimization over experiments. Odd-query errors remain given exactly by~\eqref{eq:main}; the even-query spin-flip formula alone does not assert that the same four numbers determine every odd-query error.

\section{Examples and a tensor-product obstruction}\label{sec:examples}
\subsection{Standard channels and normalization checks}
For complete depolarization, $K_{ij}=\ket i\bra j/\sqrt d$. A unit coefficient vector gives $\|K(z)\|_\HS^2=1/d$, so arithmetic--geometric mean yields $\delta=d^{-2}$, attained by $K(z)=I/d$. The exact finite error recovers
\[
b_n=\frac{q_*}{d^n m_*},
\]
and the exponent is $2\log d$. Balanced multiplicity probes and this exponential gain have precedents~\cite{chiribella}; they serve here as a check.

For qubit amplitude damping with $0\le\gamma\le1$, take
\[
K_0=\begin{pmatrix}1&0\\0&\sqrt{1-\gamma}\end{pmatrix},\qquad
K_1=\begin{pmatrix}0&\sqrt\gamma\\0&0\end{pmatrix}.
\]
Then $\det K(z)=\sqrt{1-\gamma}\,z_0^2$, so
\[
\delta=\sqrt{1-\gamma},\qquad b_{2a}=(1-\gamma)^a.
\]
For $0<\gamma<1$ the Kraus span contains no nonzero scalar multiple of a unitary: orthogonality of its columns forces the off-diagonal entry to vanish, and the two diagonal moduli then differ. Thus no positive unitary CP component exists, despite the finite positive determinant exponent.

For an orthogonal random-unitary channel $\E=\sum_jp_j\U_{U_j}$, with $\Tr(U_i^\dagger U_j)=d\delta_{ij}$,
\[
|\det K(z)|^{2/d}\le\frac{\|K(z)\|_\HS^2}{d}
=\sum_jp_j|z_j|^2\le p_{\max}.
\]
Selecting a largest-weight Kraus unitary attains equality, so $\delta=p_{\max}$. A pure replacer has matrix-space rank one and first perfect query number two.

\subsection{Werner--Holevo parity}
Consider
\[
\mathcal W_d(X)=\frac{\Tr(X)I-X^T}{d-1},\qquad
K_{ij}=\frac{\ket i\bra j-\ket j\bra i}{\sqrt{d-1}}\quad(i<j).
\]
Its Kraus span is the space of complex skew-symmetric matrices. If $d$ is odd, its maximum rank is $d-1$, so the first perfect query number is exactly $d$. If $d$ is even, an invertible skew matrix exists. Every unit coefficient vector has squared Frobenius norm $2/(d-1)$, and an equal-block skew matrix attains equal singular values. Therefore
\begin{equation}
\delta(\mathcal W_d)=\frac{2}{d(d-1)}\quad(d\text{ even}),\qquad
\xi(\mathcal W_d)=\log\frac{d(d-1)}2.
\end{equation}
At $d=2$ this is a unitary channel and the exponent correctly vanishes. At $d=3$, the first three optimal errors at $\alpha=0$ are $1,1/2,0$.

\subsection{Perfect detection need not survive a tensor product}
The next example changes the null from arbitrary local unitaries on a qutrit to arbitrary global unitaries on two qutrits. This enlarged null is essential. The underlying rank phenomenon is already known: Li, Qiao, Wang and Duan~\cite{li}, Theorem 12, show that the odd-dimensional skew space has maximum rank $d-1$ while its tensor square has full matrix rank $d^2$, using the same sum of tensor products of skew basis matrices. We use that established phenomenon and compute the exact determinant maximum for the qutrit channel, which fixes the testing exponent.

\begin{theorem}[An exact tensor-square example]\label{thm:tensor}
The channel $\mathcal W_3$ has $\delta=0$ and admits perfect detection in three queries. For the nine-dimensional channel $\mathcal W_3\otimes\mathcal W_3$, tested against arbitrary fixed $U(9)$ dynamics with the corresponding full input and output rotations,
\begin{equation}
\delta(\mathcal W_3\otimes\mathcal W_3)=\frac{2^{2/9}}{12}>0,
\qquad \xi=\log12-\frac29\log2.
\label{eq:tensor}
\end{equation}
Consequently no finite number of queries gives perfect detection in this global model. In particular, the invariant $\delta$ is not multiplicative under tensor products.
\end{theorem}
\begin{proof}
Let $A_1,A_2,A_3$ be the standard real skew basis with nonzero pairs $(1,2),(1,3),(2,3)$ and entries $+1,-1$. The Kraus matrices of the tensor square are $A_i\otimes A_j/2$. For a coefficient matrix $Z=(z_{ij})$, $\|Z\|_\HS=1$, write
\[
G(Z)=\frac12\sum_{i,j=1}^3z_{ij}A_i\otimes A_j.
\]
We claim
\begin{equation}
|\det G(Z)|=\frac{|\det Z|^3}{256}.\label{eq:tensor-det}
\end{equation}
Congruence $A\mapsto UAU^T$ acts on the three-dimensional skew space through $\bigwedge^2U$. The map $\bigwedge^2:U(3)\to U(3)$ is surjective: in a fixed exterior basis it is unitarily equivalent to $U\mapsto\det(U)\overline U$, and any target admits such a preimage by choosing a square root of its determinant. Thus arbitrary left and right unitary transformations of $Z$ correspond to unitary congruences of $G(Z)$, preserving both determinant moduli. Singular-value decomposition reduces~\eqref{eq:tensor-det} to $Z=\operatorname{diag}(s_1,s_2,s_3)$.

In this case $G(Z)$ preserves the diagonal subspace spanned by $\ket{11},\ket{22},\ket{33}$ and the three two-dimensional subspaces spanned by $\ket{ij},\ket{ji}$. Their block determinants multiply to
\[
\det G(\operatorname{diag}s)=-\frac{s_1^3s_2^3s_3^3}{256},
\]
proving the claim. Under $\sum s_i^2=1$, the maximum of $s_1s_2s_3$ is $3^{-3/2}$. Substitution into~\eqref{eq:delta} with dimension nine gives $256^{-2/9}/3=2^{2/9}/12$. The exponent and impossibility of finite perfect detection follow from Theorems~\ref{thm:exponent} and~\ref{thm:perfect}.
\end{proof}

There is no contradiction with applying a local perfect test when the null is promised to be a product unitary. A general entangling $U(9)$ does not preserve the local antisymmetric structures used by that test. The example quantifies the effect of enlarging the nuisance group, not a loss caused by freely adding a known ancillary system. The underlying matrix-space nonsingularity phenomenon was established in~\cite{li} and is related to commutative/noncommutative rank distinctions~\cite{ggow,derksen}; the exact operational exponent is the conclusion established here.


\subsection{The smallest partner: an exact depolarizing-qubit exponent}
Let $\mathcal D_2(X)=\Tr(X)I_2/2$. Theorem~\ref{thm:stable} predicts that a depolarizing qubit suffices for $\mathcal W_3$, since the latter is unital. The following calculation evaluates the exponent and proves that partner dimension two is minimal.

\begin{theorem}[Exact minimal-partner example]\label{thm:qubit-partner}
In the six-dimensional global-unitary model,
\begin{equation}\boxed{\displaystyle
\delta(\mathcal W_3\otimes\mathcal D_2)=\frac{2^{2/3}}{24},\qquad
\xi(\mathcal W_3\otimes\mathcal D_2)=\log24-\frac23\log2.}
\label{eq:qubit-partner}
\end{equation}
Thus $\mathcal W_3$ has a three-query perfect test, whereas $\mathcal W_3\otimes\mathcal D_2$ has no finite perfect test. No one-dimensional partner can produce this change.
\end{theorem}
\begin{proof}
Use the skew basis $A_1,A_2,A_3$ defined above. The twelve Kraus matrices of the product channel are $A_i\otimes E_{ab}/2$. A normalized combination is $G=H/2$, where
\[
H=\sum_{i=1}^3 A_i\otimes X_i=
\begin{pmatrix}0&X_1&X_2\\-X_1&0&X_3\\-X_2&-X_3&0\end{pmatrix},
\qquad \sum_i\|X_i\|_\HS^2=1.
\]
Let $W$ be the $3\times4$ matrix whose $i$th row lists the entries of $X_i$ in the order $(00,01,10,11)$, and let $p_{abc}$ be its minor on columns $a<b<c$, numbered $0,1,2,3$. Expanding the block determinant gives the polynomial identity
\begin{equation}
\det H=p_{013}p_{023}-p_{012}p_{123}.
\label{eq:minor-identity}
\end{equation}
This identity holds over $\CC$ without invertibility assumptions on the blocks; it can also be checked by substituting their twelve entries and collecting terms. The degree-six polynomial is, up to sign convention, the classical $2\times2\times3$ hyperdeterminant: see Lim's minors expression~\cite{lim}, slide 32, and Bremner~\cite{bremner}. Neither that invariant nor the minors expression is claimed as new here. The triangle inequality, $2|uv|\le|u|^2+|v|^2$, and Cauchy--Binet give
\[
|\det H|\le\frac12\sum_{a<b<c}|p_{abc}|^2
=\frac12\det(WW^\dagger)
\le\frac12\left(\frac{\Tr WW^\dagger}{3}\right)^3
=\frac1{54}.
\]
Equality is attained by $(X_1,X_2,X_3)=(I_2,\sigma_x,\sigma_z)/\sqrt6$. Before the factor $1/\sqrt6$, the corresponding $H$ has determinant $4$ and singular values $2,2,1,1,1,1$. Consequently
\[
\max|\det G|=\frac1{2^6\,54}=\frac1{3456}.
\]
Taking the power $2/6$ gives~\eqref{eq:qubit-partner}, and Theorem~\ref{thm:exponent} gives its exponent. Every channel on a one-dimensional system is the identity, leaving the singular skew space unchanged; hence dimension two is minimal.
\end{proof}

The stable and unstable cases now have explicit representatives with sharp constants. A pure replacer retains a finite perfect test under every partner, while the unital qutrit example loses perfect detection under the smallest possible nontrivial depolarizing partner. The blow-up criterion and the skew-space rank phenomenon are established algebra~\cite{li,degree}; the test interpretation and the two exact determinant maxima are the specific conclusions proved here.

\section{Verification, limitations and research provenance}
The proofs above are analytic. The accompanying programs provide reproducible finite checks by different expressions: direct Choi contractions in permutation projectors, Gaussian polynomial coefficient moments, normalized Schur evaluations, Takagi spectra and spin-flip spectra. The checks include 99 direct channel/query cases, 90 even-query qubit moment comparisons through twelve queries, and 252 exact rational Schur inequalities. The maximum discrepancy between direct balanced-block contractions and Gaussian moments is $2.23\times10^{-16}$. An exact symbolic calculation verifies the diagonal tensor-square determinant, and 100 complex coefficient matrices verify its general form numerically, with maximum absolute discrepancy below $2.4\times10^{-20}$. The latter checks supplement the congruence proof; they do not establish a polynomial identity on their own. For the new six-dimensional example, symbolic expansion verifies the full twelve-variable determinant identity, and 300 random complex coefficient arrays check both that identity and the sharp upper bound. An exact pure-replacer example checks the stable-class rank bound.

The general spin-flip square root was evaluated through a Hermitian positive spectral decomposition. A preliminary generic matrix-square-root computation on rank-deficient matrices produced spurious small roots; that numerical implementation was corrected before recording these results. No proof-assistant formalization is claimed.

The result requires independent unknown common input and output rotations, memoryless calls, arbitrary reference size and a fixed false-positive tolerance. It does not assert parallel optimality for calibrated noise, correlated environments, or a fixed memory budget. The determinant formula is not a channel capacity and does not imply an efficient solution of its maximization. The finite-query bounds are explicit but can be loose; a sharp polynomial prefactor for general channels remains open.

The 16 October refresh searched 909 available complete texts and 122 supplementary texts, refreshed 58 repository trees and the official source indices, and checked the current Erd\H{o}s problem index. Existing remote PDFs were not all downloaded again. The search preceding this manuscript revisited the OpenAI mathematical corpus, the author's AIRR and viXra-indexed works and research repositories, and the cited DeepMind and Anthropic research indices. The delivered search audit distinguishes automatic indexing and text triage from proof-level reading. It does not claim that every frontier publication or every proof has been examined. The proposed contribution is the universal optimizer and its determinant/rank consequences in the stated model. Worldwide novelty and research impact remain subject to expert scrutiny.

\paragraph{AI assistance.}
OpenAI Codex assisted substantially with research selection, mathematical development, writing and computation. A separate GPT-6 Astra xhigh context challenged the developing argument and contributed the signed-minor derivation and saturation check in Theorem~\ref{thm:qubit-partner}. That context therefore participated in development; its calculations are not an independent review of that result. Any additional retained review is identified by its exact manuscript hash and remains advisory and from the same provider. This is not independent human peer review. No report is attributed to Claude unless an actual report is produced and retained. The author is responsible for the submitted claims and references.

\begin{thebibliography}{99}
\bibitem{lim} L.-H. Lim, Hyperdeterminants, secant varieties, and tensor approximations, lecture slides (23 April 2008), slide 32 of 37. \href{https://www.stat.uchicago.edu/~lekheng/work/sjsu.pdf}{Author-hosted slides}.
\bibitem{bremner} M. R. Bremner, A hyperdeterminant for $2\times2\times3$ arrays, \href{https://arxiv.org/abs/1106.2988v1}{arXiv:1106.2988v1} (2011).
\bibitem{nakahira} K. Nakahira and K. Kato, Generalized quantum process discrimination problems, \emph{Phys. Rev. A} \textbf{103}, 062606 (2021); erratum \textbf{104}, 069902. \href{https://arxiv.org/abs/2104.09759}{arXiv:2104.09759}.
\bibitem{chiribella} G. Chiribella and D. Ebler, Quantum speedup in the identification of cause--effect relations, \emph{Nature Communications} \textbf{10}, 1472 (2019). \href{https://arxiv.org/abs/1806.06459}{arXiv:1806.06459}, especially Appendix B.
\bibitem{sedlak} M. Sedl\'ak and M. Ziman, Unambiguous comparison of unitary channels, \emph{Phys. Rev. A} \textbf{79}, 012303 (2009). \href{https://arxiv.org/abs/0809.4401}{arXiv:0809.4401}.
\bibitem{hayashi} M. Hayashi, Indefinite causal order strategy does not improve the estimation of group action, \emph{Quantum} \textbf{9}, 1891 (2025). \href{https://arxiv.org/abs/2501.09312}{arXiv:2501.09312}.
\bibitem{hayashi_symmetry} M. Hayashi, Y.-A. Chen, C. Zhu and X. Wang, Predicting symmetries of quantum dynamics with optimal samples, \href{https://arxiv.org/abs/2502.01464}{arXiv:2502.01464} (2025).
\bibitem{complete} L. Eriksson, A complete class of optimal quantum tests for unknown unitary dynamics, AIRR Working paper ARR-2026-5MAVBFTEF48DX8RT, v1 (16 October 2026). \href{https://airr.science/papers/ARR-2026-5MAVBFTEF48DX8RT/versions/v1/}{Permanent version}.
\bibitem{sra} S. Sra, On inequalities for normalized Schur functions, \emph{European Journal of Combinatorics} \textbf{51}, 492--494 (2016). \href{https://arxiv.org/abs/1502.04753v3}{arXiv:1502.04753v3}.
\bibitem{macdonald} I. G. Macdonald, \emph{Symmetric Functions and Hall Polynomials}, second edition, Oxford University Press (1995).
\bibitem{wootters} W. K. Wootters, Entanglement of formation of an arbitrary state of two qubits, \emph{Phys. Rev. Lett.} \textbf{80}, 2245--2248 (1998). \href{https://arxiv.org/abs/quant-ph/9709029}{arXiv:quant-ph/9709029}.
\bibitem{gour} G. Gour, Mixed-state entanglement of assistance and the generalized concurrence, \emph{Phys. Rev. A} \textbf{72}, 042318 (2005). \href{https://arxiv.org/abs/quant-ph/0506229}{arXiv:quant-ph/0506229}.
\bibitem{werner} R. F. Werner and A. S. Holevo, Counterexample to an additivity conjecture for output purity of quantum channels, \emph{J. Math. Phys.} \textbf{43}, 4353--4357 (2002). \href{https://arxiv.org/abs/quant-ph/0203003}{arXiv:quant-ph/0203003}.
\bibitem{ghosal} P. Ghosal, P. Halder, A. Patra and A. Sen(De), Quantum hypothesis testing of non-mixed-unitarity: A multifaceted hierarchy of quantum channel discrimination, \href{https://arxiv.org/abs/2609.40355v2}{arXiv:2609.40355v2} (2026).
\bibitem{ggow} A. Garg, L. Gurvits, R. Oliveira and A. Wigderson, Operator scaling: theory and applications, \href{https://arxiv.org/abs/1511.03730v4}{arXiv:1511.03730v4} (2019).
\bibitem{derksen} H. Derksen and V. Makam, On non-commutative rank and tensor rank, \href{https://arxiv.org/abs/1606.06701}{arXiv:1606.06701} (2016).
\bibitem{degree} H. Derksen and V. Makam, Polynomial degree bounds for matrix semi-invariants, \emph{Advances in Mathematics} \textbf{310}, 44--63 (2017). \href{https://arxiv.org/abs/1512.03393}{arXiv:1512.03393}, Theorem 1.8.
\bibitem{iqs} G. Ivanyos, Y. Qiao and K. V. Subrahmanyam, Constructive noncommutative rank computation is in deterministic polynomial time, \href{https://arxiv.org/abs/1512.03531v6}{arXiv:1512.03531v6} (2018).
\bibitem{li} Y. Li, Y. Qiao, X. Wang and R. Duan, Tripartite-to-bipartite entanglement transformation by stochastic local operations and classical communication and the structure of matrix spaces, \href{https://arxiv.org/abs/1612.06491v2}{arXiv:1612.06491v2} (2017), Section II.2 and Theorems 9 and 12.
\end{thebibliography}
\end{document}

